\documentclass[12pt]{article}
\usepackage[utf8]{inputenc}
\usepackage[T1]{fontenc}
\usepackage{hyperref}
\title{Southern Samo Language Resources for K-12 Curriculum Development}
\author{Compiled by Claude Sonnet 5.5 (claude-sonnet-5-5) for Andr\'as Kornai}
\date{}
\begin{document}
\maketitle

\section{Introduction}
Southern Samo (San Sud, Maka Samo; ISO 639-3 sbd; Glottolog sout2844) is an Eastern Mande language of north-western Burkina Faso, spoken mainly in Nayala, Sourou, Sanguie and Passore provinces (Nouna and Solenzo areas). Speaker estimates in the sources found vary widely (about 85,000 to 230,000). A Toma variety (also Makaa, Nyaana) is described. Jula, French and Moore are common second languages. Joshua Project reports an adult literacy rate of about 4\%. No school materials were found in search results.

\section{Primary Educational Resources}
No Southern Samo primers, textbooks or school programmes were found. A literacy centre teaching San in Toma is mentioned on the Webonary page below (the entry itself was not opened).
\begin{itemize}
\item Webonary San Sud site (introduction): \url{https://www.webonary.work/san-sud/overview/introduction/?lang=en}
\end{itemize}

\section{Linguistic and Reference Materials}
\begin{itemize}
\item Webonary San Sud dictionary entries: \url{https://www.webonary.work/san-sud/?p=9444}
\item Webonary Samo--English dictionary (C.F. Smith Jr., SIL 2018; coverage of the southern variety not verified): \url{https://www.webonary.work/samo/?p=146817}
\item Glottolog entry for Southern Samo, listing a 1985 study on South San (Toma) phonology, a 2003 study of Toma San texts (K. Woodham) and a 2001 sociolinguistic survey (J. Berthelette): \url{https://glottolog.org/resource/languoid/id/sout2844}
\item Samo language (Burkina), Wikipedia: \url{https://en.wikipedia.org/wiki/Samo_language_(Burkina)}
\item REFLEX (CNRS) bibliography page for Samo: \url{https://reflex.cnrs.fr//Lexiques/webball/biblio.php?LGNOM=Sam}
\end{itemize}

\section{Religious and Community Resources}
\begin{itemize}
\item Southern Samo in Burkina Faso, Joshua Project: \url{https://joshuaproject.net/people_groups/14717}
\end{itemize}
No Southern Samo scripture was confirmed in the searches.

\section{Digital and Audio Resources}
\begin{itemize}
\item Wikitongues page for Southern Samo (language profile; EGIDS 5 Developing): \url{https://wikitongues.org/languages/sbd/}
\end{itemize}
No audio or software resources specific to the language were found.

\section{Curriculum Development Context}
Southern Samo has a Latin-based writing system, a Webonary lexicon and some linguistic literature, but no confirmed school materials. Low literacy and the dominance of Jula and French suggest that a curriculum would need to start with orthography confirmation, reading primers and teacher training, drawing on the Webonary lexicon. The Toma and other varieties differ in intelligibility, so variety choice is an open question.

\end{document}